\section{Почему меняется квотное ранжирование}
\label{sec:ranking_shift}
Формула~\eqref{eq:selection_shift} позволяет уточнить механизм
сравнительной деградации на уровне идеальных условных целей.
Обозначим $h_i=P(D=-1\mid z_i)$ и $g_i=E[D\mid z_i]$.
Если только harm-исход получает множитель $a\ge1$, то
\begin{equation}\label{eq:harm_target}
g_{a,i}=\frac{g_i-(a-1)h_i}{1+(a-1)h_i}.
\end{equation}
Все знаменатели положительны. Приведение разности двух дробей
к общему знаменателю даёт точное тождество
\begin{equation}\label{eq:rank_numerator}
\begin{aligned}
\operatorname{sign}(g_{a,i}-g_{a,j})
&=\operatorname{sign}\bigl((g_i-g_j)+(a-1)B_{ij}\bigr),\\
B_{ij}&=h_j(1+g_i)-h_i(1+g_j).
\end{aligned}
\end{equation}
Квадратичные по $a$ произведения сокращаются. Поэтому при $g_i>g_j$
и $B_{ij}<0$ порядок меняется после единственного порога
$a_{ij}^{*}=1+(g_i-g_j)/(-B_{ij})$; при $B_{ij}\ge0$ он сохраняется.

\paragraph{Следствие для границы квоты.}
Пусть $\mathcal S$ --- исходный top-$k$ по $g$, причём все выбранные
объекты строго превосходят все невыбранные. Тогда
\begin{equation}\label{eq:critical_multiplier}
a_{\mathrm{crit}}=
\min_{\substack{i\in\mathcal S,\ j\notin\mathcal S\\B_{ij}<0}}
\left(1+\frac{g_i-g_j}{-B_{ij}}\right).
\end{equation}
Минимум по пустому множеству равен бесконечности.
До этого порога все сравнения через границу квоты сохраняют знак;
после него хотя бы один невыбранный объект превосходит выбранный.
В точке порога результат зависит от разрешения равенств.
Это доказывает условие изменения множества; величина потери accuracy
зависит также от исходного зазора полезностей.

\paragraph{Значение для выбора способа обучения.}
Порог определяется условными harm-вероятностями и зазорами
\emph{между объектами}, а не только средней ошибкой экспертов.
В приложении приведены две популяции с одной агрегированной таблицей
совместных исходов, но разной реакцией top-$k$ на harm-отбор.
Следовательно, диагностика выбора должна учитывать расположение
спасений и вреда относительно доступных признаков и границы квоты.
При этом семь идеальных gain-оценок имеют общую цель~\eqref{eq:harm_target};
различия между конечновыборочными реализациями проверяются экспериментально.
Условие выведено как алгебраическое следствие известного преобразования
при отборе, а не заявляется новым общим принципом или готовым
оценённым на данных правилом выбора. Помимо доказательства выполнены
20\,000 численных проверок попарного тождества и 250 проверок
первой смены top-$k$ на конечных популяциях.
